\section{从糟糕 BP 到整机供应链}

现在进入星海图本体。前面几章解释高继扬为什么相信 physical world AI，也解释他从自动驾驶学到了什么；这一章回答他创业时最关键的选择：为什么具身智能公司不能只做算法。

星海图的早期转折来自一个判断：如果要做具身智能，不能只做“智能”，必须做整机。高继扬在访谈中说，机器人要真正形成数据和模型闭环，必须有自己的本体、硬件、供应链和客户场景。否则基础模型再好，也缺少稳定的真实世界入口。

\begin{figure}[H]
\centering
\includegraphics[width=0.96\textwidth]{three-product-combo.png}
\caption{星海图的三件套：整机、基础模型、后训练工具。}
\end{figure}

\begin{knowledgebox}{读图：为什么是三件套}
图中三件套对应一个目标体验：像培训员工一样培训机器人。整机提供物理入口和数据采集能力；基础模型提供通用动作和理解底座；后训练工具让客户用少量示范和自我演练完成场景适配。缺少任一项，都很难把机器人从 demo 推到生产力场景。
\end{knowledgebox}

\subsection{24、25、26 的节奏}

三件套说明了最终结构，但创业公司不能一次性把所有能力铺开。本节关注节奏，因为节奏决定组织是否能承受复杂度。

高继扬给出的节奏是：2024 年重心是整机和供应链，2025 年是数据智能，2026 年开始做长期供应链。这种节奏体现了“步步为营”：长期战略坚持，但不在一个时间段内同步铺开所有事情。机器人公司如果过早同时追求硬件、模型、应用、渠道和客户，组织很容易被复杂度击穿。

\begin{importantbox}{步步为营的战略含义}
步步为营不是保守，而是按依赖关系推进。没有整机和供应链，就没有稳定数据；没有稳定数据，就没有好的 VLA；没有可用模型，就难以形成客户价值。
\end{importantbox}

\subsection{本章小结}

星海图选择整机，不是因为硬件本身更性感，而是因为整机是数据、模型和客户闭环的入口。具身智能公司要把智能做深，反而要先把物理链条做实。

